\chapter{Evaluation}
\label{ch:evaluation}

This chapter evaluates the efficacy of the TCM-Sage system compared to a non-specialized large language model (LLM) baseline. The evaluation is conducted through a blind A/B testing framework, the "Arena," involving practitioners and students from the Traditional Chinese Medicine (TCM) community.

\section{Experimental Setup}

The evaluation framework was designed to isolate the impact of the specialized RAG (Retrieval-Augmented Generation) pipeline and the TCM-specific system prompt.

\subsection{System Configurations}
Two distinct system configurations were compared:
\begin{enumerate}
    \item \textbf{Treatment (TCM-Sage):} The full pipeline described in Chapter 4, including specialized RAG context from 17 classical texts, SymMap v2.0 knowledge graph integration, and the professional TCM system prompt.
    \item \textbf{Baseline (Plain LLM):} A generic configuration using the same underlying Qwen model but with a "helpful assistant" prompt. To simulate a modern AI's capabilities, the baseline had access to \textbf{DuckDuckGo Web Search} (5 results), representing the current standard for general-purpose retrieval.
\end{enumerate}

\subsection{Blind A/B Design}
A total of 56 votes were collected from a pool of TCM students and practitioners. The evaluation used a blind design:
\begin{itemize}
    \item \textbf{Randomization:} The positions of the Treatment and Baseline responses (A/B) were randomized for each query.
    \item \textbf{Interaction:} Voters entered a clinical or theoretical query and received two simultaneous streams.
    \item \textbf{Metrics:} After reading both responses, voters could select "A is Better," "B is Better," or "Tie." Votes were stored in a JSONL format for statistical analysis.
\end{itemize}

\section{Quantitative Results}

The primary goal of the quantitative analysis was to determine if the TCM-Sage system provided a statistically significant improvement over a general-purpose LLM.

\subsection{Comparative Performance}
The distribution of the 56 votes is summarized below:
\begin{itemize}
    \item \textbf{TCM-Sage Wins:} 35 (62.5\%)
    \item \textbf{Plain LLM Wins:} 15 (26.8\%)
    \item \textbf{Ties:} 6 (10.7\%)
\end{itemize}

The TCM-Sage system outperformed the baseline in over 60\% of cases, more than doubling the win rate of the general-purpose search-augmented model.

\subsection{Statistical Significance}
To rigorously test the null hypothesis that there is no difference between the two systems, a \textbf{paired T-test} was performed on the win/loss data (assigning 1 for a win, -1 for a loss, and 0 for a tie).
\begin{itemize}
    \item \textbf{T-statistic:} 3.03
    \item \textbf{P-value:} 0.0037
\end{itemize}
With $p < 0.01$, the results are highly significant, allowing us to reject the null hypothesis.

\subsection{Effect Size}
\textbf{Cohen's d} was calculated to measure the magnitude of the improvement:
\begin{itemize}
    \item \textbf{Cohen's d:} 0.40
\end{itemize}
An effect size of 0.40 indicates a \textit{medium effect}, suggesting that while the general-purpose model is a capable baseline, the specialized TCM-Sage system provides a meaningful and consistent advantage in the quality of clinical reasoning and source accuracy.

\section{Temporal Analysis}

The evaluation was conducted over a four-day period (April 3 to April 6, 2026), during which the system underwent minor iterative refinements based on initial feedback.

\begin{itemize}
    \item \textbf{April 3:} 2 TCM-Sage / 1 Plain (67\% Win Rate) — Initial testing phase.
    \item \textbf{April 4:} 13 TCM-Sage / 8 Plain (62\% Win Rate) — Main testing batch from student groups.
    \item \textbf{April 5:} 8 TCM-Sage / 0 Plain (100\% Win Rate) — Testing following the integration of the \textit{Measurement Conversion} table.
    \item \textbf{April 6:} 12 TCM-Sage / 6 Plain (67\% Win Rate) — Extended testing with a wider group of practitioners.
\end{itemize}

The cumulative trend shows a consistent preference for TCM-Sage, with a notable peak in performance after specific clinical knowledge (e.g., unit conversions) was hard-coded into the system prompt.

\section{Qualitative Feedback and Case Studies}

Qualitative feedback from TCM professionals provided deeper insights into the specific strengths and weaknesses of both systems.

\subsection{Dosage and Unit Conversion}
A doctoral student at the Guangdong Provincial Hospital of TCM tested a query regarding the dosage of \textit{Mahuang Tang}. He observed that the general LLM (Baseline) correctly identified modern dosage ranges but the TCM-Sage system (at that time) provided a literal translation of Han Dynasty units that the model incorrectly converted (1 \textit{liang} $\approx$ 3g). 

\textit{Outcome:} This feedback led to the immediate implementation of the \textit{Measurement Conversion} table (1 \textit{liang} $\approx$ 13.8g) in the system prompt. Subsequent tests showed that practitioners found the corrected system far superior to the baseline, as it provided both historical accuracy and modern relevance.

\subsection{Source Traceability}
Students consistently highlighted the \textbf{Citation Panel} as the most valuable feature. One student noted: "The ability to see the original classical text immediately, without having to copy-paste the formula name into a search engine to verify the source, is the key value of this system." This indicates that the \textit{traceability} provided by RAG is as important as the \textit{content} of the generated response.

\subsection{Diagnostic Hedging and Safety}
A senior practitioner noted that the early version of the system was "over-assertive" in its diagnosis of \textit{Pi Xu} (Spleen Deficiency) based on limited symptoms. Following this, \textbf{Hedging Principles} were added to the prompt, requiring the system to use more tentative language ("The symptoms suggest a possibility of...") when data is incomplete. The practitioner later confirmed that the updated responses felt more like a professional consultation than a "black-box" diagnosis.

\subsection{Novel Clinical Insights}
In one recorded case, a student consulted the system for a sore throat. The system suggested \textit{Jiegeng Tang}, a classical formula the student had not considered. Upon reviewing the provided citations, the student confirmed the validity of the recommendation. This demonstrates the system's potential as a "clinical co-pilot" that can surface relevant but overlooked classical knowledge.

\section{Limitations of Evaluation}

While the results are statistically significant, several limitations must be acknowledged:
\begin{enumerate}
    \item \textbf{Sample Size:} 56 votes are sufficient for statistical significance but do not represent a comprehensive clinical trial.
    \item \textbf{Voter Demographic:} The majority of voters were TCM students rather than senior clinicians with decades of experience.
    \item \textbf{Evaluation Format:} The Arena tests overall user preference, which includes factors like formatting and tone, not just clinical accuracy.
    \item \textbf{Prompt Influence:} Because the system prompts differed between the two models, the evaluation measures the effectiveness of the \textit{entire pipeline} (Prompt + RAG + KG), not just the retrieval component in isolation.
\end{enumerate}

\section{Discussion of Results}

The results of the evaluation demonstrate that a specialized RAG system tailored for TCM significantly outperforms a general-purpose search-augmented LLM.

\begin{itemize}
    \item \textbf{Accuracy vs. Verifiability:} The primary advantage of TCM-Sage lies in its ability to provide verbatim classical citations. While the general LLM can often "summarize" TCM knowledge from the web, it frequently halluncinates formula origins or misinterprets classical syntax.
    \item \textbf{The Value of Domain Knowledge:} The integration of a specialized knowledge graph and a professional system prompt (including ancient-to-modern measurement conversions) is essential for clinical utility. Without these, even a high-quality RAG system is prone to fundamental errors in interpretation.
    \item \textbf{Impact on the Field:} The $p = 0.0037$ result strongly suggests that specialized AI tools can provide a higher standard of information for the TCM community than general-purpose assistants, particularly in the areas of citation traceability and classical methodology adherence.
\end{itemize}
